% typography and font
% ====================================================================
\usepackage[utf8]{inputenc}				% to be able to type accented characters
\usepackage[T1]{fontenc}				% so that LaTeX is able to manage accented characters
% \usepackage{lmodern}					% font choice
\usepackage[french, english]{babel}		% document typographic rules
\usepackage[usenames,dvipsnames,svgnames,table]{xcolor}	% text color
\usepackage{csquotes}					% Context sensitive quotation capabilities

% figures
% ====================================================================
\usepackage{graphicx}				% NB: to be commented if standalone package is loaded
\usepackage{float}					% imporve floating object management, add positioning option H
\usepackage[font=small]{caption}	% caption options
\usepackage{placeins}				% usage: \FloatBarrier
\usepackage{tikz}					% drawing library
\usetikzlibrary{shapes,arrows,shadows,patterns,pgfplots.groupplots,decorations.markings,chains,calc}


% layout
% ====================================================================
\usepackage[top=30mm,bottom=30mm]{geometry}		% page dimension use the showframe option to see the different areas of the page
\usepackage{fancyhdr}		% controls headers and footers
\usepackage{minitoc}		% table of content for each chapter
\usepackage{titlesec}		% customisation of section macros (titles, headers and contents)
\usepackage{epigraph}		% chapter quote
\usepackage{authblk} 		% footnote style & author/affiliation
\usepackage{pdfpages}		% insert pdf file
\setcounter{secnumdepth}{3}	% section numbering depth


% bibliography options
% ===============================================
\usepackage[
style = numeric-comp, % define both the bibliography style and the citation style
% refsection = chapter, % This option automatically starts a new reference section at a document division such as a chapter or a section
% bibstyle = ... to specify the bibliography style (see documentation)
% citestyle = .... to specify the citation style (see documentation)
sorting=none, % sorting option (here by name then year then title)
% subentry,
sortcites = true,% Whether or not to sort citations if multiple entry keys are passed to a citation command
maxnames = 2, % if more authors than maxnames the list of authors is truncated to "minnames"
minnames = 2, 
doi=false,isbn=false, % disable printing for selected fields
% maxbibnames = 3,
hyperref=true, %to move back and forth between the main matter and the bibliography (needs the hyperref package to be loaded, see below)
refsegment=chapter %alt : section
]{biblatex}

% To be used for subbibliography at the end of each chapter (optional)
\defbibheading{subbibliography}{%
  \section*{References for Chapter \ref{refsegment:\therefsection\therefsegment}}}

% math
% ====================================================================
\usepackage{amsmath}		% math environments
\usepackage{amssymb}		% math symbols
\usepackage{bbold}			% usage: so that \mathbb{1} give an indicator function
\usepackage{bm,bbm}			% bold in math mode
\usepackage{cancel}			% usage: strike through terms in equation 
\usepackage[normalem]{ulem}	% usage: strikethrough horizontaly in text mode, normalem so that \emph{} remains unchanged
\usepackage{empheq}			% usage: box in math aligned environnemnt, left brace, ... 
\usepackage{marvosym}
\usepackage{manfnt}
\usepackage{mathrsfs}
% \usepackage{unicode-math}

% tabular
% ====================================================================
% \usepackage{longtable}  % usage: table than span over several pages
% \usepackage{tabularx}   % usage: [among others] tables with automatic line wrap
\usepackage{ltablex}	  % combine <tabularx> and <longtable>
% == tabular options
\setlength{\tabcolsep}{6pt}			% default value - space between colums in tabular
\renewcommand{\arraystretch}{1.2} 	% default value - space between lines in tabular
% == end tabular options


% Colors (official Inria colors)
% ========================================================================
\definecolor{myskyblue}{RGB}{81,167,249}
\definecolor{myred}{RGB}{230,51,18} 
\definecolor{mygrayblue}{RGB}{56,66,87} 
\definecolor{myyellow}{RGB}{255,205,28} 
\definecolor{mylightgreen}{RGB}{199,214,79} 
\definecolor{mymauve}{RGB}{101,97,169}
\definecolor{myorange}{RGB}{240,126,38} 
\definecolor{mygreen}{RGB}{149,193,31} 
\definecolor{myblue}{RGB}{20,136,202} 
\definecolor{mylilas}{RGB}{155,0,79} 

% theorem definition
% ====================================================================
\newtheorem{theorem}{Theorem}[section]
\newtheorem{corollary}{Corollary}[theorem]
\newtheorem{lemma}[theorem]{Lemma}


% others
% ====================================================================
\usepackage{fontawesome}		% icon bank
\usepackage{multirow}			% multi rows in table
\usepackage{array}				% allows to set array column width
\usepackage{siunitx}			% manage numbers and units
% == siunitx options
\sisetup{detect-weight = true}		% siunitx option, put text in bold if inside a bold environment
\sisetup{range-phrase=--}			% for \SIrange{1}{10}{\second} replace "1s to 10s" by "1s -10 s"
\sisetup{range-units=single}		% for \SIrange{1}{10}{\second} replace "1s to 10s" by "1-10 s"
\DeclareSIUnit\molar{\textsc{M}}	% define <molar> unit
% == end siunitx options
\usepackage[version = 4]{mhchem}	% chemical equations
\usepackage{xcolor}					% usage: colorbox



\extrafloats{50} % number of floats that can be processed by LaTeX

% hyperlinks
% ====================================================================
\usepackage[bookmarks,
			plainpages=false,
			pdfpagelabels,
			%% pdfpagemode=FullScreen,  
			colorlinks=true,
			linkcolor=orange,
			anchorcolor=gray,
			citecolor=gray,
			urlcolor=gray,
			hyperindex=true,
			breaklinks]
			{hyperref}
